\chapter{Conclusiones provisionales}
\label{ch:conclusions}

\textbf{La fase se mejoró mediante una corrección espectral parcialmente
predecible, aunque no se reconstruyó de forma completamente autónoma.}
El resultado central de S4 es distinguir una componente de retardo compartida
y transferible, una dirección común al arreglo condicionada a la captura y un
residual relativo con estructura espectral. Esta distinción permite interpretar
qué se ha logrado y elegir el tratamiento adecuado de lo que permanece abierto.


La primera etapa de investigación establece un marco reproducible de comparación entre
un canal medido y su representación electromagnética, y distingue fidelidad macroscópica,
fidelidad de magnitud, concordancia de fase y respuesta a perturbaciones físicas. La
cadena desde escena hasta observable permite separar calibración del entorno, tratamiento
instrumental y aprendizaje del problema inverso.

Sobre DICHASUS dc01, la calibración B0--B2 mejora potencia, magnitud y dispersión de
retardos, sin una mejora correspondiente de fase bruta. La comparación se sitúa en la
línea de calibración diferenciable de materiales y no identifica por sí sola parámetros
constitutivos verdaderos. Los resultados justifican evaluar expresamente qué propiedades
no quedan determinadas por el objetivo macroscópico de entrenamiento.

S4.2 caracteriza una corrección afín de fase formada por retardo efectivo e interceptos
por enlace. El retardo es relativo al modelo y a la referencia de procesamiento; no se
identifica automáticamente con STO físico. La transferencia entre antenas y frecuencias
muestra estructura compartida, con diferencias claras entre interpolación espectral y
extrapolación entre medias bandas. La alineación oráculo conserva su función diagnóstica
y no se interpreta como predicción autónoma.

S4.3 documenta estructura espectral residual y asociaciones con la distribución energética
de las trayectorias. El análisis entre posiciones, dentro de posición y mediante
remuestreo agrupado evidencia que el nivel de agregación modifica la interpretación.
La asociación observada no determina una causa única ni reemplaza un contraste explícito
de independencia que conserve la distribución marginal y el procedimiento de ajuste.

S4.4 muestra que la partición original favorece interpolación densa: la distancia mediana
al vecino de entrenamiento es 0.0483 m. La evaluación posterior mediante cinco bloques
espaciales y exclusión de 0.5 m produce un MAE medio de retardo de 27.56 ns con posición,
19.29 ns con descriptores RT y 18.69 ns con ambos. La información del trazado conserva
utilidad en regiones reservadas de dc01, aunque la mejora depende de la región y no
constituye todavía transferencia entre escenas.

S4.5 vincula el retardo predicho con métricas del canal. En interpolación densa,
el error angular mediano disminuye de 81.42 a 45.59 grados y la coherencia
aumenta de 0.156 a 0.634. Bajo evaluación por bloques,
el predictor combinado alcanza error angular mediano de 64.08 grados y coherencia
mediana de 0.448, frente a 81.49 grados y 0.157 con corrección exclusiva de CPO.
El CPO continúa ajustándose con la medición evaluada: la contribución demostrada es
la transferencia del retardo y su utilidad condicionada, no la recuperación completa
de fase sin acceso al canal. La degradación respecto de interpolación densa cuantifica
la importancia del protocolo espacial.

S4.6 amplía el diagnóstico a 640\,000 fases residuales en 10\,000 capturas.
La resultante global de 0.001088 y las medianas por antena de 0.01024 y 0.00630
no respaldan una fase absoluta fija fuertemente concentrada. La dirección común
estimada con las 32 antenas reduce el error mediano de 88.99 a 45.97 grados;
esa reducción es descriptiva y requiere controles de ajuste finito para
cuantificar su significación. La concentración marginal baja no excluye por
sí sola sesgos relativos estables entre antenas.

La recuperación temporal produjo correspondencia XYZ exacta y un manifiesto
sin tiempos faltantes. Con esa referencia, copiar la fase común del vecino
espacial alcanza errores medianos de 89.75 y 88.85 grados; copiarla del vecino
temporal produce 90.34 y 90.56 grados. La falta de una ventaja clara se observa
también en los intervalos temporales examinados. El resultado delimita el
alcance de esas reglas locales, sin demostrar independencia, equivalencia
estadística o una causa instrumental única.

La etapa siguiente S4.7a evaluará concordancia relativa respecto de fase común.
Los productos espaciales $G$ eliminan además el retardo común; por ello su
contribución predictiva debe medirse con productos entre frecuencias $K$ o con
la distancia del canal apilado respecto de una sola rotación por arreglo.
S4.7b comparará incertidumbre circular y por trayectoria cuando el residual lo justifique.
S5 evaluará representaciones con entradas, invariancias y costes declarados;
S6 contrastará tanto $\Delta H$ como la sensibilidad del observable invariante
mediante desplazamientos conocidos y calibración basal fija; S7 validará
operadores Wi-Fi y LoRa específicos. La recuperación respiratoria y la
incertidumbre de la tarea permanecen hipótesis por validar con mediciones
independientes. Una invariancia que mejora fidelidad estática puede eliminar
parte de la sensibilidad física y, por tanto, debe evaluarse frente a la tarea.

El criterio doctoral de éxito consiste en establecer qué representación y qué tratamiento
de la discrepancia conservan información útil para sensado, bajo qué condiciones esa
información se transfiere y cuándo el sistema debe reconocer que la evidencia disponible
es insuficiente. La comparación entre error de reconstrucción y error diferencial
proporciona el vínculo matemático y experimental entre los avances de canal y ese objetivo.

La propuesta actualizada convierte ese criterio en dos ramas complementarias y una
secuencia de intervenciones físicas. La rama robusta buscará estabilidad frente a
transformaciones instrumentales especificadas; la coherente conservará micromovimiento
cuando exista una referencia verificable. Su separación constituye una hipótesis a
contrastar, porque una transformación común puede contener simultáneamente deriva y
señal física. Por ello la invariancia estática y la sensibilidad diferencial se medirán
juntas, sin convertir una mejora de normalización en evidencia automática de sensado.

La geometría se tratará en tres escalas: recinto, registro de antenas y desplazamiento
local. G0--G4 permitirá medir el aporte de reconstrucción y calibración; E0--E2 establecerá
el piso instrumental y la respuesta mecánica antes de E3--E4 con personas. E5 evaluará
transferencia real entre condiciones independientes. El cronograma y la propuesta de
Ciencia de Frontera priorizan referencia RF, mecánica y trazabilidad, reutilizando el
equipo reportado. Los resultados actuales terminan en S4.6d.2; el programa restante
expresa compromisos de validación y decisiones experimentales, no logros anticipados.
